Of 16 comparisons where validation selected a non-zero weight, \textbf{6 (37.5\%) were worse than
\texttt{w\ =\ 0}} on the evaluation block. Binomial \(p = 0.45\) against a coin flip; 95\% CI
{[}0.15, 0.65{]}. We initially reported this as evidence that the selection signal is noise.

It is not. At n = 16 the design detects a proportion only below \textbf{0.147} or above \textbf{0.853}.
Power against the observed effect is \textbf{9.5\%}. Reaching 80\% power would need \textbf{125
comparisons}. A selection procedure with a genuinely useful 30\% error rate would have been
missed 91\% of the time.

The correct statement is \emph{undetermined}, not \emph{chance}. We also retract our earlier framing of
four configurations as independent replications: they share assets and periods.

\textbf{And ``undetermined'' is still not the whole answer.} \citet{cawley2010over} established
that a model-selection criterion is itself an estimator with a variance, and that optimising it
over a finite sample over-fits the criterion exactly as training over-fits the data. Their
analysis applies to this design directly, and sharpens the verdict from ``underpowered'' to
\textbf{the comparison was conducted inside the noise floor of its own selection step}. \Cref{app:selection}
gives the argument, the quotations, and the remedy we did not run.

\emph{Control: a power calculation. Cost: free. Never performed until a reviewer demanded it.
The relevant literature is from 2010 and we did not read it until after the result.}
